\chapter{Multi-Curve and Collateral Discounting}\label{ch:rc:multi-curve-and-collateral-discounting}

Until the summer of 2007 a three-month interbank loan and a three-month loan
against collateral cost almost the same: the spread between them was about ten
basis points, and in euros under six. From August 2007 to May 2009 it stayed
above fifty basis points, and after Lehman Brothers failed it was above a
hundred for six months. Every swap pricer on the Street rested on one
assumption that this broke: that a single curve both projects the floating
rates a swap will pay and discounts every cash flow, so that a floating leg on
the three-month rate is worth par. Once three-month and six-month interbank
rates, and the overnight rate, each carried their own premium for bank credit
and liquidity, there was no single curve left. Dealers moved in a few years to
a framework with one curve per index for projection and one curve per
collateral agreement for discounting, and in 2020 the clearing houses moved
their discounting for the whole market in two weekends. This chapter builds
that framework on chapter~1's curve builder.

\section{The 2007 break and the tenor basis}

Before 2007 the textbook argument of One Quant Book 2, chapter 9, held: a
floating leg paying the rate for each period, on a curve that both projects
and discounts, is worth $P(t_0)-P(t_n)$. If a bank can borrow for six months
at six-month Euribor and roll it, a six-month floating leg is worth par and
three- and six-month legs are worth the same. After 2007 a lender for six
months demanded more than two three-month loans in a row, because of the
chance that the borrower fails in the second three months and the value of
keeping one's cash available. The difference is paid in a swap of two
floating legs.

\begin{definition}[Tenor basis, tenor basis swap]\label{def:rc:multi-curve-and-collateral-discounting:basis}
A \emph{tenor basis swap}\index{tenor basis swap} exchanges two floating legs
on indices of different tenors in the same currency, for example six-month
Euribor against the overnight rate compounded, with a spread on one leg that
makes it worth zero at inception. The \emph{tenor basis}\index{tenor basis} at
a maturity is that par spread.
\end{definition}

\begin{proposition}[Tenor basis from two swap curves]\label{prop:rc:multi-curve-and-collateral-discounting:basisrates}
If an overnight-index swap and an interbank-index swap of the same maturity
have identical fixed legs and are discounted on the same curve, the \omterm{def:rc:multi-curve-and-collateral-discounting:basis}{tenor
basis} spread (paid on the overnight leg, on the fixed leg's schedule) equals
the difference of their par rates, $b_n = S^{\text{IRS}}_n - S^{\text{OIS}}_n$.
\end{proposition}
\begin{proof}
Each par swap equates its floating leg with $S_n A$, with the same annuity $A$.
Receiving the interbank leg and paying the overnight leg plus $b$ is worth
$S^{\text{IRS}}_nA - S^{\text{OIS}}_nA - bA$, zero for $b=S^{\text{IRS}}_n-S^{\text{OIS}}_n$.
\end{proof}

\section{Discount curves and projection curves}

\begin{definition}[Discount curve, projection curve, multi-curve framework]\label{def:rc:multi-curve-and-collateral-discounting:curves}
A \emph{discount curve}\index{discount curve} gives the present value of a
cash flow paid under a given collateral agreement; a \emph{projection
curve}\index{projection curve} of an index gives its forward fixings, as
$F(t;T_1,T_2) = \bigl(P^{\text{proj}}(t,T_1)/P^{\text{proj}}(t,T_2)-1\bigr)/\delta$,
without being used to discount anything. The \emph{multi-curve
framework}\index{multi-curve framework} values a swap by projecting each
floating fixing on its index's projection curve and discounting every cash
flow on the discount curve of the trade's collateral.
\end{definition}

The \omterm{def:rc:multi-curve-and-collateral-discounting:curves}{projection curve}'s ``discount factors'' are only a device for storing
forward fixings; in the language of One Quant Book 4, chapter 5, the forward
of each fixing is its expectation under the forward measure of the
\emph{discount} curve, which is why the \omterm{def:rc:multi-curve-and-collateral-discounting:curves}{projection curve} is calibrated with the
\omterm{def:rc:multi-curve-and-collateral-discounting:curves}{discount curve} held fixed.

\begin{method}[Two-stage calibration]\label{met:rc:multi-curve-and-collateral-discounting:twostage}
(1) Calibrate the \omterm{def:rc:multi-curve-and-collateral-discounting:curves}{discount curve} from overnight-index instruments
(\cref{met:rc:curve-construction:newton}). (2) Holding it fixed, calibrate the
\omterm{def:rc:multi-curve-and-collateral-discounting:curves}{projection curve} of each index from that index's fixing, futures and swaps,
each priced as $\sum_j F_j\delta_jP^{\text{disc}}(t_j)$ against $K\sum_i
\delta_iP^{\text{disc}}(t_i)$. (3) Check that every instrument reprices and
that the forward basis (projected fixing minus the overnight forward for the
same period) is positive and smooth.
\end{method}

\omcode{code/firm/multicurve/firm_multicurve.py}{93}{115}{Calibration of a projection curve on a fixed discount curve.}\label{lst:rc:multi-curve-and-collateral-discounting:proj}

\begin{example}[A euro market]\label{ex:rc:multi-curve-and-collateral-discounting:euro}
Take illustrative €STR swap rates from 1.95\% at one year to 2.55\% at ten
and 2.60\% at thirty, and six-month Euribor swap rates 20 basis points higher
at one year, 15 at ten and 12 at thirty, with a six-month fixing of 2.13\%. The
\omterm{def:rc:multi-curve-and-collateral-discounting:curves}{projection curve} reprices all ten Euribor instruments; by
\cref{prop:rc:multi-curve-and-collateral-discounting:basisrates} the \omterm{def:rc:multi-curve-and-collateral-discounting:basis}{tenor basis} is
20, 15 and 12 basis points at one, ten and thirty years. The forward basis,
fixing by fixing (\cref{fig:rc:multi-curve-and-collateral-discounting:forwards}), is
21.6 basis points on the first six-month period and 11.8 on the period starting
in ten years: a par basis is an average of forward bases.
\end{example}

\begin{omfigure}
\begin{tikzpicture}
\begin{axis}[width=12cm,height=5.6cm,
  xlabel={start of the six-month period (years)},ylabel={forward rate (\%)},
  tick label style={font=\small},label style={font=\small},
  xmin=0,xmax=30,ymin=1.8,ymax=3.4,grid=major,grid style={gray!25},
  legend columns=2,legend style={font=\footnotesize,at={(0.5,-0.25)},anchor=north,draw=none,fill=none,/tikz/every even column/.append style={column sep=8pt}},
  table/col sep=comma]
\addplot[omThm,very thick] table[x=t,y=euribor]{figdata/rates-credit-risk/02-multi-curve-and-collateral-discounting/forwards.csv};
\addplot[omDef,thick,dashed] table[x=t,y=ois]{figdata/rates-credit-risk/02-multi-curve-and-collateral-discounting/forwards.csv};
\legend{six-month Euribor forward (projection curve),six-month €STR forward (discount curve)}
\end{axis}
\end{tikzpicture}

\omcaption{Forward rates for six-month periods from the two euro curves of
\cref{ex:rc:multi-curve-and-collateral-discounting:euro}. The gap is the
forward \omterm{def:rc:multi-curve-and-collateral-discounting:basis}{tenor basis}: about 22 basis points on the first period, 12 after ten
years, 9 at the long end. Data: the chapter's illustrative curves and
tutorial.}\label{fig:rc:multi-curve-and-collateral-discounting:forwards}
\end{omfigure}

\begin{example}[One swap, two frameworks]\label{ex:rc:multi-curve-and-collateral-discounting:compare}
A receiver of 3.50\% against six-month Euribor for ten years on EUR~100
million, with the ten-year Euribor swap at 2.70\%, is worth
EUR~7\,135\,408 in the \omterm{def:rc:multi-curve-and-collateral-discounting:curves}{multi-curve framework} and EUR~7\,075\,108 on a single
curve built from the Euribor swaps, which discounts at the higher Euribor
rates. Both frameworks agree on the par rate, 2.70\%, because each is
calibrated to it; they disagree by EUR~60\,300 on an off-market swap.
\end{example}

\section{Collateral discounting and the credit support annex}

Why discount at the overnight rate? Because that is the rate paid on the cash
that secures the trade. Under a credit support annex (One Quant Book 2,
chapter 10), a party whose derivative has positive value $V_t$ receives
$V_t$ in cash from the other, and pays the agreed interest on it; the
derivative is, in effect, funded by its own collateral.

\begin{definition}[Collateral rate, CSA discounting]\label{def:rc:multi-curve-and-collateral-discounting:csa}
The \emph{collateral rate}\index{collateral rate} $c_t$ of a trade is the rate
the collateral taker pays on cash collateral under the trade's collateral
agreement, typically the overnight benchmark of the collateral currency.
\emph{CSA discounting}\index{CSA discounting} values each cash flow of a fully
cash-collateralised trade on the curve of its collateral rate.
\end{definition}

\begin{proposition}[The collateral rate discounts]\label{prop:rc:multi-curve-and-collateral-discounting:csa}
If a trade is continuously and fully collateralised in cash paying $c_t$, and
the hedges it needs are financed at the risk-free rate, its value is
\[
V_t = \E^{\mathbb Q}_t\Bigl[\exp\Bigl(-\int_t^T c_u\,du\Bigr)V_T\Bigr].
\]
\end{proposition}
\begin{proof}[Partial proof]
The holder of the trade holds $V_t$ in collateral, owed back, on which it pays
$c_t$; the hedging positions earn and cost $r_t$ and cancel. The account
$V_t$ therefore grows at $c_t$: under $\mathbb Q$, $dV_t = c_tV_t\,dt +
dM_t$ with $M$ a martingale, whose solution is the formula. A full proof
(Piterbarg, 2010) keeps track of the funding of the hedges.
\end{proof}

\begin{omfigure}
\begin{tikzpicture}[font=\small,
  box/.style={draw,thick,rounded corners=2pt,align=center,minimum height=1.1cm,minimum width=3.2cm,inner sep=4pt},
  arr/.style={-{Stealth[length=2.5mm]},thick}]
\node[box,draw=omDef,fill=omDef!8] (bk) at (0,0) {bank\\(derivative worth $V_t>0$)};
\node[box,draw=omThm,fill=omThm!8] (cp) at (7.6,0) {counterparty\\(owes $V_t$)};
\draw[arr] (bk.15) -- node[above=6pt] {swap cash flows} (cp.165);
\draw[arr] (cp.195) -- node[below=6pt] {swap cash flows} (bk.345);
\draw[arr,omProp] (cp.north) to[bend right=25] node[above=3pt] {cash collateral $V_t$ (variation margin)} (bk.north);
\draw[arr,omExo] (bk.south) to[bend right=25] node[below=3pt] {interest $c_t V_t\,dt$ at the collateral rate} (cp.south);
\end{tikzpicture}

\omcaption{A cash-collateralised derivative. The bank holds collateral equal
to the trade's value and pays the \omterm{def:rc:multi-curve-and-collateral-discounting:csa}{collateral rate} on it; its position in the
trade is funded by the collateral, so its value accrues at $c_t$ and is
discounted on the curve of $c_t$
(\cref{prop:rc:multi-curve-and-collateral-discounting:csa}).}\label{fig:rc:multi-curve-and-collateral-discounting:csa}
\end{omfigure}

A cleared swap is the same, with the clearing house as counterparty: variation
margin moves every day, and interest on it flows the other way.

\begin{definition}[Price alignment interest]\label{def:rc:multi-curve-and-collateral-discounting:pai}
\emph{Price alignment interest}\index{price alignment interest} (PAI) is the
interest a clearing house charges the receiver of cash variation margin and
pays to the payer, at an overnight rate it chooses, so that a cleared swap is
economically equivalent to a bilateral swap collateralised in cash at that
rate. The clearing house discounts at the PAI rate.
\end{definition}

\section{Collateral choice and cross-currency collateral}

Many collateral agreements let the poster choose among several currencies or
securities, and substitute them later.

\begin{definition}[Collateral choice option]\label{def:rc:multi-curve-and-collateral-discounting:choice}
A \emph{collateral choice option}\index{collateral choice option} is the
right of the collateral poster to choose, and later change, which of several
eligible collaterals it posts. Expressed in the trade's currency, the poster
chooses the collateral earning the highest rate, so the trade is discounted at
$\max_i c^{(i)}_t$ over the eligible collaterals.
\end{definition}

Collateral in a foreign currency earns that currency's overnight rate; swapped
into the trade currency at the cross-currency basis (One Quant Book 2,
chapter 16), it earns the domestic overnight rate plus or minus the basis.
With the basis for lending euros against dollars at $b(t)$, posting euros
earns $r_{\text{SOFR}}-b$ in dollar terms: more than dollar cash when $b<0$.

\begin{example}[A dollar trade with a two-currency CSA]\label{ex:rc:multi-curve-and-collateral-discounting:choice}
Take an illustrative basis of $-15$ basis points at the front rising linearly
to $+5$ at ten years. Euros are the better collateral until 7.5 years, dollars
after. A payment of USD~100 million in ten years is worth USD~68\,598\,292
under a dollar-only CSA and USD~68\,213\,510 when the payer may post either
currency: the choice, used at every date, raises the effective discount rate
by an average of $56.25/10 = 5.6$ basis points a year and costs the receiver
USD~384\,782. With volatile rates and basis the option is worth more than this
intrinsic value.
\end{example}

\begin{omfigure}
\begin{tikzpicture}
\begin{axis}[width=12cm,height=5.4cm,
  xlabel={time (years)},ylabel={rate in dollars (\%)},
  tick label style={font=\small},label style={font=\small},
  xmin=0,xmax=10,ymin=3.1,ymax=4.4,grid=major,grid style={gray!25},
  legend columns=3,legend style={font=\footnotesize,at={(0.5,-0.25)},anchor=north,draw=none,fill=none,/tikz/every even column/.append style={column sep=8pt}},
  table/col sep=comma]
\addplot[omDef,thick] table[x=t,y=usd]{figdata/rates-credit-risk/02-multi-curve-and-collateral-discounting/collateral.csv};
\addplot[omProp,thick,dashed] table[x=t,y=eur]{figdata/rates-credit-risk/02-multi-curve-and-collateral-discounting/collateral.csv};
\addplot[omThm,line width=2.2pt,opacity=0.45] table[x=t,y=eff]{figdata/rates-credit-risk/02-multi-curve-and-collateral-discounting/collateral.csv};
\legend{dollar cash (SOFR),euro cash swapped to dollars,effective (the maximum)}
\draw[gray,dashed] (axis cs:7.5,3.1) -- (axis cs:7.5,4.4);
\end{axis}
\end{tikzpicture}

\omcaption{\omterm{def:rc:curve-construction:pillar}{Instantaneous forward rates} earned by two eligible collaterals, in
dollars, and the effective discount rate when the poster may choose. The
lines cross at 7.5 years (dashed), where the illustrative basis changes sign.
Data: the chapter's tutorial.}\label{fig:rc:multi-curve-and-collateral-discounting:choice}
\end{omfigure}

\section{After the interbank rates: the 2020 discounting switches}

When overnight benchmarks replaced EONIA and the dollar effective federal funds
rate as the reference rates of their markets, the clearing houses changed the
rate of their \omterm{def:rc:multi-curve-and-collateral-discounting:pai}{price alignment interest}, and with it the \omterm{def:rc:multi-curve-and-collateral-discounting:curves}{discount curve} of every
cleared swap in the currency.

\begin{definition}[Discounting switch]\label{def:rc:multi-curve-and-collateral-discounting:switch}
A \emph{discounting switch}\index{discounting switch} is a change of the
\omterm{def:rc:multi-curve-and-collateral-discounting:csa}{collateral rate} of a population of trades, and hence of their \omterm{def:rc:multi-curve-and-collateral-discounting:curves}{discount curve},
on a set date, with the resulting changes in value compensated in cash and,
where required, the changes in risk compensated with basis swaps.
\end{definition}

\begin{dated}{2026-09}{The 2020 switches}\label{dat:rc:multi-curve-and-collateral-discounting:switches}
EONIA was computed as €STR plus 8.5 basis points from 1 October 2019 until it
was discontinued on 3 January 2022. LCH moved discounting and price alignment
on all euro swaps from EONIA to €STR flat on 27 July 2020, compensating each
account in cash for the change in net present value, with projected cash flows
held fixed (a ``constant-forward'' method). It moved dollar swaps from federal
funds to SOFR over the weekend of 16--19 October 2020 with cash for the value
change and federal-funds/SOFR basis swaps for the risk change: more than one
million contracts and USD~120 trillion of notional; clients who did not want
the basis swaps had them auctioned, USD~24 billion net notional, to 18 dealers
at close to zero cost.
\end{dated}

Because EONIA was fixed at €STR plus 8.5 basis points, the euro switch was a
deterministic shift of the \omterm{def:rc:multi-curve-and-collateral-discounting:curves}{discount curve}: every discount factor rose by
$e^{0.00085\,t}$. An in-the-money swap, whose future cash flows are net
receipts, gained value; an out-of-the-money swap lost; the compensation
reversed each change.

\omcode{code/rates-credit-risk/02-multi-curve-and-collateral-discounting/python/rc_multicurve.py}{125}{135}{The clearing house's constant-forward valuation of the switch.}\label{lst:rc:multi-curve-and-collateral-discounting:switch}

\begin{omfigure}
\begin{tikzpicture}
\begin{axis}[width=11.5cm,height=5.2cm,ybar,bar width=14pt,
  ylabel={value change (EUR thousand)},xlabel={remaining maturity},
  tick label style={font=\small},label style={font=\small},
  xmin=-0.6,xmax=7.6,ymin=0,ymax=650,grid=major,grid style={gray!25},
  xtick=data,xticklabels from table={figdata/rates-credit-risk/02-multi-curve-and-collateral-discounting/switch.csv}{years},
  nodes near coords,nodes near coords style={font=\scriptsize},
  table/col sep=comma]
\addplot[fill=omDef!55,draw=omDef] table[x=i,y=change]{figdata/rates-credit-risk/02-multi-curve-and-collateral-discounting/switch.csv};
\end{axis}
\end{tikzpicture}

\omcaption{Gain at the euro \omterm{def:rc:multi-curve-and-collateral-discounting:switch}{discounting switch} of a receiver of 1.50\% against
six-month Euribor on EUR~100 million, by remaining maturity, on illustrative
curves of July 2020 (euro swap rates between $-0.50\%$ and $+0.08\%$). The
clearing house's cash compensation took each gain back. Data: the chapter's
tutorial.}\label{fig:rc:multi-curve-and-collateral-discounting:switch}
\end{omfigure}

\section{Tutorial: a euro market on two curves}\label{tut:rc:multi-curve-and-collateral-discounting:tutorial}

\begin{tutorial}
\textbf{Goal.} Build an €STR \omterm{def:rc:multi-curve-and-collateral-discounting:curves}{discount curve} and a six-month Euribor \omterm{def:rc:multi-curve-and-collateral-discounting:curves}{projection
curve}, read the \omterm{def:rc:multi-curve-and-collateral-discounting:basis}{tenor basis}, value a swap in both frameworks, and value the
2020 \omterm{def:rc:multi-curve-and-collateral-discounting:switch}{discounting switch}. \textbf{End state:}
\cref{fig:rc:multi-curve-and-collateral-discounting:forwards,fig:rc:multi-curve-and-collateral-discounting:switch}
and the numbers of \cref{ex:rc:multi-curve-and-collateral-discounting:compare}.
\begin{tutsteps}
\item \textbf{\omterm{def:rc:multi-curve-and-collateral-discounting:curves}{Discount curve}}: \texttt{discount()}, nine €STR swaps with
chapter~1's builder, monotone convex.
\item \textbf{\omterm{def:rc:multi-curve-and-collateral-discounting:curves}{Projection curve}}: \texttt{projection()} calibrates the fixing and
nine Euribor swaps with the \omterm{def:rc:multi-curve-and-collateral-discounting:curves}{discount curve} held (\cref{lst:rc:multi-curve-and-collateral-discounting:proj}); check
\texttt{basis\_table()} returns 20, 19, 18, 17, 16, 15, 14, 13 and 12 basis
points.
\item \textbf{Two frameworks}: \texttt{discounting\_comparison()}.
\item \textbf{The switch}: \texttt{switch\_value(20, 0.015, 5e8)};
\texttt{fig\_rc\_multicurve.py} writes the charts.
\end{tutsteps}
\textbf{What to change next.} Make the basis flat at 15 basis points and check
that the forward basis becomes flat too; value the switch with the \omterm{def:rc:multi-curve-and-collateral-discounting:curves}{projection
curve} recalibrated on the €STR \omterm{def:rc:multi-curve-and-collateral-discounting:curves}{discount curve} instead of held fixed, and
compare with the clearing house's method.
\end{tutorial}

\section{Build: discount and projection curves}\label{bld:rc:multi-curve-and-collateral-discounting:multicurve}

\begin{build}
\bfield{Purpose} Every rates instrument of the book is priced with a \omterm{def:rc:multi-curve-and-collateral-discounting:curves}{projection
curve} for its index and a \omterm{def:rc:multi-curve-and-collateral-discounting:curves}{discount curve} for its collateral.
\bfield{Interface} \texttt{IborSwap}, \texttt{Fixing};
\texttt{calibrate\_projection(spot, disc, instruments, kind)};
\texttt{ibor\_swap\_pv(proj, disc, \ldots)}; \texttt{tenor\_basis};
\texttt{ShiftedCurve(base, shift)}; \texttt{CollateralChoiceCurve(base,
spreads)}; all curves expose \texttt{df(date)}.
\bfield{Rules} The \omterm{def:rc:multi-curve-and-collateral-discounting:curves}{discount curve} is calibrated first and never changed by the
projection calibration; fixed legs annual ACT/360, floating legs on the
index's tenor; collateral choice as the pointwise maximum of eligible rates;
\texttt{firm\_curvebuild} and Book 2's \texttt{firm\_curve} imported, not edited.
\bfield{Acceptance tests} \texttt{code/firm/multicurve/tests/}: the \omterm{def:rc:multi-curve-and-collateral-discounting:curves}{projection
curve} reprices its swaps on its \omterm{def:rc:multi-curve-and-collateral-discounting:curves}{discount curve}; the \omterm{def:rc:multi-curve-and-collateral-discounting:basis}{tenor basis} equals the
par-rate difference; with no basis, projection and discount forwards agree;
payer and receiver values are opposite; a collateral choice never raises the
value of a receivable and reduces to the single curve with one collateral.
\bfield{Stretch} Several indices (one, three, six, twelve months) from basis
swaps; a cross-currency \omterm{def:rc:multi-curve-and-collateral-discounting:curves}{discount curve} from cross-currency basis swaps; the
collateral choice with volatile basis by simulation.
\end{build}

\begin{omsources}
\item V.\ Piterbarg, ``Funding beyond discounting: collateral agreements and
derivatives pricing'', \emph{Risk}, February 2010; ``Cooking with collateral'',
\emph{Risk}, 2012.
\item A.\ De Socio, ``The interbank market after the financial turmoil'', Banca
d'Italia, Temi di Discussione 819, 2011.
\item LCH, \emph{SwapClear transition to €STR discounting}; \emph{SOFR
discounting: plan for the SwapClear compensation process}, 2019--2020.
\item M.\ Henrard, \emph{Interest Rate Modelling in the Multi-Curve Framework},
Palgrave Macmillan, 2014.
\end{omsources}

\section{Exercises}

\begin{exercise}[$\star$]\label{exo:rc:multi-curve-and-collateral-discounting:1}
The ten-year swap against six-month Euribor is at 2.70\% and the ten-year
€STR swap at 2.55\%, with identical annual fixed legs. What is the ten-year
\omterm{def:rc:multi-curve-and-collateral-discounting:basis}{tenor basis}, and on which leg is it paid?
\end{exercise}

\begin{exercise}[$\star$]\label{exo:rc:multi-curve-and-collateral-discounting:2}
Before 2007 the three-month unsecured-secured spread was about ten basis
points; afterwards above fifty. Explain in two sentences why that ended the
single-curve framework.
\end{exercise}

\begin{exercise}[$\star$]\label{exo:rc:multi-curve-and-collateral-discounting:3}
A euro swap is collateralised in euro cash paying €STR. Which curve discounts
it, which projects its six-month fixings, and what changes if the CSA pays
€STR plus 10 basis points?
\end{exercise}

\begin{exercise}[$\star\star$]\label{exo:rc:multi-curve-and-collateral-discounting:4}
In \cref{ex:rc:multi-curve-and-collateral-discounting:choice}, compute the
integral of the effective spread over ten years and the resulting discount
factor ratio, and check the USD~384\,782.
\end{exercise}

\begin{exercise}[$\star\star$]\label{exo:rc:multi-curve-and-collateral-discounting:5}
Explain the EUR~60\,300 of
\cref{ex:rc:multi-curve-and-collateral-discounting:compare}: which way does the
single curve err, and why are both frameworks right about the par rate?
\end{exercise}

\begin{exercise}[$\star\star$]\label{exo:rc:multi-curve-and-collateral-discounting:6}
At the euro switch, a client is \emph{paying} 1.50\% on EUR~100 million for
ten years. Using \cref{fig:rc:multi-curve-and-collateral-discounting:switch},
give its value change and who pays the compensation to whom.
\end{exercise}

\begin{exercise}[$\star\star\star$]\label{exo:rc:multi-curve-and-collateral-discounting:7}
\emph{Coding.} With \texttt{forward\_table()}, give the forward basis on the
six-month period starting in ten years, compare with the ten-year par basis of
15 basis points, and explain the difference.
\end{exercise}

\begin{exercise}[$\star\star\star$]\label{exo:rc:multi-curve-and-collateral-discounting:8}
\emph{Find the flaw.} ``Our CSA lets the counterparty post euros or dollars;
today dollars are cheaper to post, so the option is out of the money and we
can value the trade on the dollar curve.'' Correct it.
\end{exercise}

\section{Problem: The Switch Weekend}

\begin{problem}[{Weekend problem --- a pension fund's swap at the euro discounting switch}]\label{pb:rc:multi-curve-and-collateral-discounting:1}
On Friday 24 July 2020 a pension fund receives 1.50\% annually against
six-month Euribor on EUR~500 million, cleared at LCH, with twenty years left;
it was traded in 2010. On the illustrative July 2020 curves of the chapter
(€STR swaps from $-0.50\%$ to $-0.05\%$, Euribor swaps 12 to 17 basis points
higher), the twenty-year Euribor par rate is 0.08\%. On Monday 27 July the
clearing house discounts at €STR flat instead of EONIA ($=$ €STR $+$ 8.5
basis points).

\medskip\noindent\textbf{Part I --- The position.}
\begin{enumerate}
\item Is the swap an asset or a liability of the fund, and why?
\item Give its value on Friday under EONIA discounting.
\item Which cash flows of the swap are known, which are projected?
\item What does the fund receive on its variation margin, before and after the switch?
\item Why does the discount rate of a cleared swap equal the PAI rate?
\end{enumerate}

\medskip\noindent\textbf{Part II --- The switch.}
\begin{enumerate}[resume]
\item By how much does every discount factor to $t$ years change?
\item Give the swap's value on Monday by the clearing house's constant-forward method.
\item Give the value change. Who gains?
\item Who pays the cash compensation, how much, and to whom?
\item Why hold the projected cash flows fixed rather than recalibrate the
\omterm{def:rc:multi-curve-and-collateral-discounting:curves}{projection curve} on the new \omterm{def:rc:multi-curve-and-collateral-discounting:curves}{discount curve}?
\end{enumerate}

\medskip\noindent\textbf{Part III --- The risk.}
\begin{enumerate}[resume]
\item After the switch, how much does the swap lose if the €STR \omterm{def:rc:multi-curve-and-collateral-discounting:curves}{discount curve}
rises by one basis point, projections held?
\item Before the switch the same number described EONIA risk. What risk has
changed, and why did the euro switch need no basis-swap compensation while the
dollar switch did?
\item Would a payer of 1.50\% have gained or lost at the switch?
\item Why is the thirty-year receiver's gain more than twice the fifteen-year's
(\cref{fig:rc:multi-curve-and-collateral-discounting:switch})?
\item Does the switch change the fund's economic position? Its accounts?
\end{enumerate}

\medskip\noindent\textbf{Part IV --- Judgement.}
\begin{enumerate}[resume]
\item Why did the market move all cleared swaps on one weekend rather than
each trade at its own pace?
\item What would a bilateral trade under a CSA paying EONIA have needed?
\item Could a dealer have profited from knowing the switch date? Why not much?
\item State the \emph{named result}: the fund's value change and compensation,
and its discount sensitivity after the switch.
\item In one sentence: what decides the \omterm{def:rc:multi-curve-and-collateral-discounting:curves}{discount curve} of a derivative?
\end{enumerate}
\end{problem}

\section{Interview questions}

\begin{interviewq}[$\star$ \iqroles{researcher, bank}]\label{iq:rc:multi-curve-and-collateral-discounting:1}
Why do we discount collateralised swaps at the overnight rate and not at the
index they pay?
\end{interviewq}

\begin{interviewq}[$\star\star$ \iqroles{researcher, developer}]\label{iq:rc:multi-curve-and-collateral-discounting:2}
Walk me through building a six-month Euribor curve after 2008. What is
calibrated first, and why?
\end{interviewq}

\begin{interviewq}[$\star\star$ \iqroles{trader}]\label{iq:rc:multi-curve-and-collateral-discounting:3}
What is a \omterm{def:rc:multi-curve-and-collateral-discounting:basis}{tenor basis swap}, and why is the six-month basis positive?
\end{interviewq}

\begin{interviewq}[$\star\star$ \iqroles{bank, risk}]\label{iq:rc:multi-curve-and-collateral-discounting:4}
A CSA allows collateral in three currencies. How does that change the value of
the trades under it, and which way?
\end{interviewq}

\begin{interviewq}[$\star\star\star$ \iqroles{researcher, risk}]\label{iq:rc:multi-curve-and-collateral-discounting:5}
A clearing house changes its discounting rate. What happens to the value and
the risk of a swap book, and how would you compensate members fairly?
\end{interviewq}

\begin{interviewq}[$\star\star\star$ \iqroles{developer}]\label{iq:rc:multi-curve-and-collateral-discounting:6}
Your pricer uses the \omterm{def:rc:multi-curve-and-collateral-discounting:curves}{projection curve} to discount by mistake. Which trades
show the error, and how large is it?
\end{interviewq}
